\section*{الفصل \ref{ch:b2:alkene-redox} --- أكسدة الألكينات واختزالها}

\begin{solution}{exo:b2:alkene-redox:1}
ميثيل هكسان حلقي؛ 1-ميثيل هكسان حلقي-1-ول؛ trans-2-ميثيل هكسان حلقي-1-ول (إضافة متزامنة لذرتي H
وB، فتكون \ce{OH} والهيدروجين الجديد في الجهة نفسها، والميثيل في وضع trans بالنسبة إلى \ce{OH})؛
أكسيد 1-ميثيل هكسين حلقي؛ cis-1-ميثيل هكسان حلقي-1،2-ديول؛ 6-أوكسو هبتانال
\ce{CH3CO(CH2)4CHO}.
\end{solution}

\begin{solution}{exo:b2:alkene-redox:2}
الهيدروبورة--الأكسدة: 2-ميثيل بروبان-1-ول. والإماهة المحفَّزة بالحمض:
2-ميثيل بروبان-2-ول.
\end{solution}

\begin{solution}{exo:b2:alkene-redox:3}
trans-2،3-ثنائي ميثيل أوكسيران، ومجموعتا الميثيل على جهتين متقابلتين من الحلقة. وهو كيرالي (لا
مستوي مرآة فيه) ويتكوّن خليطًا راسيميًا.
\end{solution}

\begin{solution}{exo:b2:alkene-redox:4}
المعالجة اللاحقة المختزلة: جزيئان من البروبانال. والمعالجة اللاحقة المؤكسدة: جزيئان من حمض
البروبانويك.
\end{solution}

\begin{solution}{exo:b2:alkene-redox:5}
(أ) ميزو-بوتان-2،3-ديول، غير نشط لأنه لاكيرالي. (ب) الديولان
(2\textit{R}،3\textit{R}) و(2\textit{S}،3\textit{S}) بكميتين متساويتين، غير نشط
لأنه راسيمي.
\end{solution}

\begin{solution}{exo:b2:alkene-redox:6}
في الحمض يهاجم الماء الكربون الثالثي، وفي القاعدة يهاجم الهيدروكسيد \ce{CH2}؛ وفي
الحالتين تنتهي مجموعة \ce{OH} على كل كربون: فيتكوّن 2-ميثيل بروبان-1،2-ديول نفسه. ومع الميثانول يعطي
الطريقان إيثرين مختلفين (كما في الفصل).
\end{solution}

\begin{solution}{exo:b2:alkene-redox:7}
الهدرجة بمكافئ واحد من \ce{H2} على حفاز بلاديوم مسموم
(حفاز ليندلار): فالإضافة المتزامنة إلى الرابطة الثلاثية تتوقف عند الألكين (\textit{Z}).
\end{solution}

\begin{solution}{exo:b2:alkene-redox:8}
يتصل البروبانون (3 كربونات) والبيوتانال (4 كربونات) عند كربونيهما الكربونيليين: \ce{(CH3)2C=CH-CH2CH2CH3}،
أي 2-ميثيل هكس-2-ين.
\end{solution}

\begin{solution}{exo:b2:alkene-redox:9}
\ce{BH3} (أو بوران ضخم)، ثم \ce{H2O2} و\ce{NaOH}: الهكسان-1-ول. أما الإماهة المحفَّزة بالحمض
فتمر بالكاتيون الكربوني الثانوي وتعطي الهكسان-2-ول (ماركوفنيكوف)، مع
نواتج ثانوية معاد ترتيبها.
\end{solution}

\begin{solution}{exo:b2:alkene-redox:10}
\omterm{def:b2:alkene-redox:epoxide}{حمض البيروكسي}، وهو محب للإلكترونات، يتفاعل أسرع مع الرابطة المزدوجة الحلقية الأكثر استبدالًا والأغنى
بالإلكترونات. والبوران الضخم يتفاعل مع \ce{C=CH2} الطرفية الأقل إعاقة.
\end{solution}

\begin{solution}{exo:b2:alkene-redox:11}
trans: \omterm{def:b2:alkene-redox:epoxide}{حمض بيروكسي}، ثم الماء في الحمض (فتح مضاد). cis: \ce{OsO4} بكمية حفزية مع
عامل معيد للأكسدة، أو برمنغنات مخففة باردة (إضافة متزامنة).
\end{solution}

\begin{solution}{exo:b2:alkene-redox:12}
للمركب \ce{C6H10} درجتا عدم تشبع؛ ومكافئ واحد من \ce{H2} يعني رابطة \ce{C=C} واحدة، فثمة
حلقة واحدة. والثنائي الألدهيد الوحيد ذو الكربونات الستة يأتي من حلقة انفتحت عند رابطتها المزدوجة:
الهكسين الحلقي.
\end{solution}

\begin{solution}{pb:b2:alkene-redox:1}
\textbf{1.} $10 \times 12.011 + 16 \times 1.008 = \qty{136.24}{g/mol}$.
\textbf{2.} $(2 \times 10 + 2 - 16)/2 = 3$.
\textbf{3.} حلقة واحدة ورابطتان مزدوجتان.
\textbf{4.} \ce{C10H16 + 2H2 -> C10H20}، أي 1-إيزوبروبيل-4-ميثيل هكسان حلقي.
\textbf{5.} $1.00/136.24 = \qty{7.34}{mmol}$ من الليمونين، و\qty{14.68}{mmol} من \ce{H2}.
\textbf{6.} $V = nRT/p = 0.01468 \times 8.314 \times 298.15/10^5 = \qty{3.64e-4}{m^3}$،
أي \qty{364}{mL}.
\textbf{7.} لا: فكل من كربوني الحلقة 1 و4 يحمل مسارين حلقيين متماثلين؛ والمتماكبان cis 
وtrans لاكيراليان.
\textbf{8.} الميثانال \ce{HCHO} (كربون واحد) من الطرف \ce{=CH2}، ومركب واحد ذو تسع كربونات \ce{C9}،
هو 3-أسيتيل-6-أوكسو هبتانال \ce{CH3CO-CH2CH2-CH(COCH3)-CH2-CHO}.
\textbf{9.} كربوناها ينتميان إلى الحلقة نفسها: فكسر الرابطة المزدوجة يفتح الحلقة
دون أن يفصل شيئًا.
\textbf{10.} الميثانال (غاز).
\textbf{11.} تصير مجموعات الألدهيد أحماضًا كربوكسيلية (ويستمر الميثانال حتى \ce{CO2})؛ 
والكيتونات لا تتغير.
\textbf{12.} الكيتون والألدهيد في السلسلة \ce{C9} كانا متصلين في الحلقة؛ والكيتون الثاني
(\ce{COCH3}) والميثانال كانا طرفي الرابطة المزدوجة في السلسلة الجانبية.
\textbf{13.} مركز واحد (الكربون C4 في الحلقة)؛ ولا يمسه \omterm{def:b2:alkene-redox:cleavage}{التحلل الأوزوني} فيبقى
مركزًا فراغيًا في الناتج \ce{C9}.
\textbf{14.} الرابطة \ce{C=CH2} الخارج حلقية، الأقل إعاقة من الرابطة المزدوجة الحلقية الثلاثية الاستبدال.
\textbf{15.} 2-(4-ميثيل هكس-3-ين حلقي-1-يل)بروبان-1-ول: \ce{CH2OH} عند الكربون الطرفي
السابق.
\textbf{16.} أولي.
\textbf{17.} تذهب \ce{OH} إلى الكربون الأقل استبدالًا، عكس التوجيه الذي
تتنبأ به قاعدة ماركوفنيكوف للإماهة.
\textbf{18.} ترتيبان فراغيان عند المركز الجديد، يعطيان متماكبين لامرآويين (مع المركز
الموجود سابقًا).
\textbf{19.} الإماهة المحفَّزة بالحمض (أو إماهة أخرى وفق ماركوفنيكوف): الكحول الثالثي،
$\alpha$-تربينيول.
\textbf{20.} الرابطة المزدوجة الحلقية: ثلاثية الاستبدال، وأغنى بالإلكترونات.
\textbf{21.} 1،2-إيبوكسي-1-ميثيل-4-(بروب-1-ين-2-يل)هكسان حلقي، والأكسجين يجسر بين C1 وC2.
\textbf{22.} عند C1، الكربون الأكثر استبدالًا (الحامل للميثيل).
\textbf{23.} 1-ميثيل-4-(بروب-1-ين-2-يل)هكسان حلقي-1،2-ديول، ومجموعتا \ce{OH} في وضع trans (فتح
مضاد).
\textbf{24.} \omterm{def:b2:alkene-redox:dihydroxylation}{الهدركسلة الثنائية} المتزامنة للرابطة المزدوجة الحلقية (\ce{OsO4} بكمية حفزية)، التي تعطي الديول
cis --- إذا أمكن جعل الكاشف يتفاعل مع تلك الرابطة المزدوجة انتقائيًا.
\textbf{25.} يمتص \qty{1.00}{g} من الليمونين $\boldsymbol{\approx \qty{364}{mL}}$ من
ثنائي الهيدروجين.
\end{solution}
